\documentclass[11pt,a4paper]{amsart}

\usepackage{cgeometry}

\addbibresource{ComplexGeometry.bib}
\defbibheading{bibliography}[\refname]{\section*{#1}}
\setlength{\bibitemsep}{0pt}
\renewcommand*{\bibfont}{\footnotesize}

\allowdisplaybreaks[2]
\emergencystretch=1em

%%%%%%%%%%%%%%%%%%%%%%%%%%%%%%%%%%%%%%%%%%%%%%%%%%%%%%%%%%%%%%%%%%%%%%%%%%%%%
%% Notation
%%%%%%%%%%%%%%%%%%%%%%%%%%%%%%%%%%%%%%%%%%%%%%%%%%%%%%%%%%%%%%%%%%%%%%%%%%%%%

\title[Second-stage affine star]{A second-stage obstruction on a nonambient affine tropical star}
\author{}
\date{24 September 2026}

\begin{document}

\begin{abstract}
We give a compact affine $G$-tropical space whose Chambert--Loir--Ducros plurisubharmonic hierarchy does not collapse after its first decreasing-limit stage. A bounded continuous second-stage function has no pointwise approximation by ordinary smooth plurisubharmonic functions on any fixed neighborhood of the central node. The example uses a restricted rational affinage on an infinite star. It is a counterexample in the affine $G$-tropical category; no Berkovich analytic realization of its restricted smooth structure is asserted.
\end{abstract}

\maketitle

\section{Introduction}

The transfinite plurisubharmonic sheaf of Chambert--Loir and Ducros starts with ordinary smooth functions with positive smooth Hessian, then repeatedly takes local decreasing sequential limits \cite[Definition~8.1.1 and Lemma~8.1.4]{CLD12}. Its first stage consists of functions locally approachable by one decreasing sequence of smooth plurisubharmonic functions on a fixed ordinary neighborhood \cite[Definition~8.1.13]{CLD12}. We show that first-stage collapse fails for an affine structure obtained by restricting the permitted slopes on the compact infinite star of \cite[\S6.7.3]{CLD12}.

\begin{maintheorem}\label{thm:second-stage-star}
There is a compact connected affine $G$-tropical space $X$ and a bounded continuous function $u\in\operatorname{Psh}_{X,2}(X)$ whose germ at a point $o$ is not in $\operatorname{Psh}_{X,1}$. More strongly, no sequence of ordinary smooth CLD plurisubharmonic functions on any fixed ordinary neighborhood of $o$ converges pointwise to $u$ there.
\end{maintheorem}

This settles the unrestricted first-stage assertion for affine $G$-tropical spaces negatively. The construction does not supply an analytic space with this affinage. Hence it does not decide the version requiring a Berkovich analytic counterexample, for which compatibility with analytic logarithmic coordinates and scalar extension would also be necessary.

\section{A restricted affinage on the compact star}\label{sec:affinage}

We use the rational parlinear convention of \cite[\S\S1.8.2 and~1.8.8]{CLD12}. Let $D$ be the countable branch set
\[
 D=\{r_j,b_j:j\geq1\}\cup\{a_{j,n}:j,n\geq1\}.
\]
Let $X$ be the compact star of \cite[\S6.7.3]{CLD12} with branch set $D$: each branch has coordinate $t\in[-\infty,0]$, all points $t=0$ form the node $o$, and every ordinary neighborhood of $o$ contains all but finitely many complete branches. The domain category and the standard parlinear sheaf are those of that example. Write $e_d$ for the coordinate vector of branch $d$, and put
\begin{equation}\label{eq:slope-space}
 L_{\mathbb Q}=\operatorname{span}_{\mathbb Q}\{\ell_{j,n}:j,n\geq1\},\qquad
 \ell_{j,n}=e_{r_j}-n e_{b_j}-e_{a_{j,n}}.
\end{equation}
Every vector in this span has finite support. Set $L_{\mathbb R}=L_{\mathbb Q}\otimes_{\mathbb Q}\mathbb R$.

On a domain $W$ of the standard star, let $A(W)$ consist of continuous functions that are rational affine on each connected branch portion, constant near every included $-\infty$ endpoint, and whose slopes on terminal branch intervals present at $o\in W$ extend to some vector in $L_{\mathbb Q}$. There is no node-slope condition when $o\notin W$; separated branch components away from the node have their ordinary rational affine structure. Let $P$ be the subsheaf of the standard parlinear sheaf whose sections are $G$-locally sections of $A$.

\begin{lemma}\label{lem:site-inheritance}
The inherited domain category with parlinear sheaf $P$ is a rational $G$-tropical space, and $A$ is an affinage. On an ordinary open neighborhood of $o$, every $A$-coordinate has a unique full terminal slope vector in $L_{\mathbb Q}$.
\end{lemma}

\begin{proof}
The slope-extension condition is stable under restriction: discarding a terminal branch condition preserves the same extension vector. It is also ordinary local on each domain. A relative ordinary neighborhood of $o$ in $W$ retains a terminal subinterval of every branch present at $o$; on each connected terminal portion, a rational affine function has one slope. Thus one local extension vector already gives the required slopes throughout those portions. Other components impose no node condition. Constants belong to $A$, and $A$ is a rational vector subpresheaf of the standard parlinear sheaf.

The $G$-local definition makes $P$ a $G$-sheaf. Every finite $P$-map is a map for the standard star, so the inherited domain, inverse-image, and compact-image axioms continue to hold. To check composition with a parlinear target map $\phi$, take its $G$-cover by rational affine cells. Their preimages under the given map are inherited domains and form a $G$-cover: for each source point, finitely many target cells through its image cover an ordinary target neighborhood, whose inverse image is an ordinary source neighborhood inside their finite preimage union. On each preimage, $\phi$ is a rational affine combination of the components and hence lies in $P$; the $G$-sheaf property glues these expressions. These are the axioms and site-morphism condition of \cite[\S\S3.1.2--3.1.4]{CLD12}. The ordinary locality of $A$ and its $G$-generation of $P$ are precisely the affinage conditions \cite[\S\S6.1.1--6.1.2]{CLD12}.

Finally, an ordinary open node neighborhood has a terminal interval in every branch. Its actual slope vector is therefore uniquely determined, and the defining extension condition puts it in $L_{\mathbb Q}$.
\end{proof}

Ordinary smoothness at $o$ means a smooth real function of finitely many $A$-coordinates on one ordinary neighborhood \cite[Definition~6.2.1]{CLD12}. Differentiating such a presentation at $o$ shows that every ordinary smooth function has terminal slope vector in $L_{\mathbb R}$. On a finite branch interior the affine structure is the usual one-dimensional rational structure, so positivity of the smooth Hessian is ordinary convexity. At an included $-\infty$ endpoint every $A$-coordinate is locally constant. Consequently each ordinary smooth psh function is convex along branches and constant near their included endpoints. We will use this necessary condition, not its converse at arbitrary node germs.

\section{The second-stage function}\label{sec:construction}

Choose a smooth convex nondecreasing $\rho:\mathbb R\to[0,\infty)$ that is zero on $(-\infty,-1]$ and has derivative $1$ on $[-1/2,\infty)$. Put $c=\rho(0)>0$. For each $j$ define $h_j$ on $X$ by $h_j=\rho(t)-c$ on $b_j$ and $h_j=0$ elsewhere. Thus $-c\leq h_j\leq0$ and $h_j(o)=0$.

\begin{lemma}\label{lem:first-stage-blocks}
Each $h_j$ belongs to $\operatorname{Psh}_{X,1}(X)$.
\end{lemma}

\begin{proof}
Remove the closed tail $\{r_j\}\times[-\infty,-1]$ and call the resulting ordinary open node neighborhood $B_j$. For $n\geq1$ define $f_{j,n}$ on $B_j$, with value $2c/n$ at $o$, by
\[
 f_{j,n}=\begin{cases}
 \rho(t)-c+2c/n,&\text{on }b_j,\\
 -t/n+2c/n,&\text{on }r_j,\\
 (\rho(t)-c)/n+2c/n,&\text{on }a_{j,n},\\
 2c/n,&\text{on all other branches.}
 \end{cases}
\]
Near $o$ it is affine with slope vector $-\ell_{j,n}/n\in L_{\mathbb Q}$, so it is itself an $A$-coordinate there and has zero smooth Hessian. On branch interiors the formulas are smooth and convex; at included $-\infty$ endpoints they are locally constant. The only transition Hessians are $\rho''(t)$ and $\rho''(t)/n$, both nonnegative. Positivity is $G$-local \cite[Definition~3.5.1]{CLD12}, so $f_{j,n}$ is ordinary smooth CLD psh on $B_j$.

The four branch formulas give, everywhere on the fixed $B_j$,
\begin{equation}\label{eq:block-error}
 \frac cn\leq f_{j,n}-h_j\leq\frac{2c+1}{n}.
\end{equation}
Choose increasing integers $n_k$ with $n_{k+1}\geq((2c+1)/c)n_k$. Then \eqref{eq:block-error} implies $f_{j,n_{k+1}}\leq f_{j,n_k}$, and the decreasing limit is $h_j$ on $B_j$. At every point away from $o$, the function $h_j$ is locally ordinary smooth psh. Ordinary locality in the first stage therefore gives $h_j\in\operatorname{Psh}_{X,1}(X)$.
\end{proof}

Put $u_m=\sum_{j=1}^{m}2^{-j}h_j$. Finite sums and positive scalings preserve the first stage: at each point intersect finitely many approximation neighborhoods and add their decreasing smooth psh rows termwise. Since $h_j\leq0$, the sequence $u_m$ decreases. It converges uniformly to
\begin{equation}\label{eq:witness}
 u=\sum_{j\geq1}2^{-j}h_j,
\end{equation}
because $\|2^{-j}h_j\|_\infty\leq c2^{-j}$. Hence $u$ is bounded and continuous and belongs to $\operatorname{Psh}_{X,2}(X)$ by the successor-stage definition \cite[Definition~8.1.1]{CLD12}.

\begin{proof}[Proof of \Cref{thm:second-stage-star}]
It remains to exclude first-stage approximation at $o$. Let $V$ be any ordinary neighborhood of $o$. The cofinite node topology gives an index $j$ for which all branches $r_j,b_j,a_{j,n}$, $n\geq1$, lie wholly in $V$. Let $f$ be any ordinary smooth CLD psh function on $V$. Its terminal slope vector $s$ lies in $L_{\mathbb R}$ by the smooth-presentation rule in \cref{sec:affinage}. Along every full branch of block $j$, convexity and constancy near $-\infty$ imply that its terminal slope is nonnegative.

The block-$j$ component of $s$ is a finite sum $\sum_n\alpha_n\ell_{j,n}$. On each $a_{j,n}$ its slope is $-\alpha_n\geq0$, while on $r_j$ the slope is $\sum_n\alpha_n\geq0$. Since only finitely many $\alpha_n$ are nonzero, all vanish. In particular the terminal slope on $b_j$ is zero. A convex function whose derivative is zero near $-\infty$ and at $0$ is constant throughout that branch; hence $f|_{b_j}=f(o)$. The same equality holds for every pointwise limit of ordinary smooth psh functions on $V$.

But \eqref{eq:witness} gives $u(o)=0$ and $u(b_j,-\infty)=-2^{-j}c$. Thus $u$ is no pointwise limit of such functions on $V$. As $V$ was arbitrary, its germ is not in the first stage.
\end{proof}

\printbibliography

\end{document}
